\section{Análisis de Incidentes y Accidentes (Casos Reales)}

\textbf{Objetivo General.}-- Motivar mediante la exposición de incidentes y accidentes reales, el análisis de las situaciones para buscar con calma y razonamiento lógico la solución que dé lugar a la exposición mínima, evitando así acciones precipitadas que pueden conducir a riesgos mayores.

\subsection{Procedimiento de Emergencia}

El alumno explicará cómo proceder en cada uno de los tipos de incidentes más comunes, en los que esté a su alcance la posibilidad de solución. Indicará a quién acudir cuando no disponga de los medios para resolver el caso, señalando las medidas o disposiciones inmediatas que debe tomar para evitar que se agrave la situación.

\subsection{Informe de Incidentes y Accidentes}

El alumno indicará los tipos de incidentes que deberá informar a sus superiores y la forma en que lo hará deberá estar consciente de que todo accidente se informará de inmediato.

\subsection{Plan de Emergencia}

\textbf{Práctica No. 8.--} Manejo de un equipo radiográfico, solución de fallas. (2 hr)

\textbf{Práctica No. 9.--} Recuperación de fuentes, uso de manipuladores a distancia (2h)

\textbf{Práctica No. 10.--} Simulacro de un accidente. (4 hr)

\textbf{Evaluación final teórico-práctica.}

Duración total del curso: 40 Hrs. (24 hrs teoría, 16 hrs práctica)